\SourceLeaf{157}
Fährt man auf dieselbe Weise fort, so erhält $\Omega$ die Gestalt
\[
\Omega = (\nu C)^2+(\nu' C')^2+(\nu'' C'')^2+\cdots+(\nu^{(n-1)} C^{(n-1)})^2+(\beta^{(n)} B^{(n)})^2,
\]
wobei alle $C$ von $x$ unabhängig sind.



\Needspace{6\baselineskip}
\begin{center}
\textsc{Gauss} an \textsc{Olbers}. Göttingen, 14. April 1819.
\end{center}

\begin{center}
. . . . . . . . . . . . . . . .
\end{center}

Ebenso ist mir Ihr Einwand gegen die richtige Formel für das Gewicht des Ergebnisses aus einer Kombination von 2 Bestimmungen mit den Gewichten $p$, $q$, nämlich $\dfrac{pq}{p+q}$, nicht recht verständlich, es sei denn, Sie haben meinen Brief nicht mit der Stelle bei \textsc{Littrow} verglichen. Allerdings setze ich voraus, dass die Teilergebnisse voneinander \textls{unabhängig} sind. Es kommt aber darauf an, von \emph{welcher Art} die Kombination ist. Im \emph{vorliegenden Fall} ist der Umstand, dass $\dfrac{pq}{p+q}$ stets $<p$ und $<q$ ist, nicht widersinnig. Das Gegenteil wäre widersinnig. Nehmen wir einen anderen ganz ähnlichen Fall. \textsc{Schumacher} beabsichtigt, einen Längengradbogen zu messen. Seine äußersten Punkte seien $A$ und $C$. Er kann sie, das heißt seine Raketen, nicht unmittelbar vergleichen und braucht daher einen Zwischenpunkt $B$. Vergleicht er zuerst $A$ mit $B$ 100-mal und danach $C$ mit $B$ ebenfalls 100-mal, so sind die Teildifferenzen der geografischen Länge $B-A$, $C-B$ in hohem und gleichem Maße zuverlässig. Ihr Gewicht wird durch $100=p$, $100=q$ angegeben. Der aus der Verbindung beider abgeleitete Unterschied zwischen den Orten $C$, $A$ wird aber offenbar weniger zuverlässig sein, als wenn er aus 100 unmittelbaren Vergleichen abgeleitet worden wäre. Um wie viel geringer das Gewicht des Ergebnisses aus dieser Kombination ist, kann nur die Wahrscheinlichkeitstheorie lehren. Sie ergibt, dass dieses \textls{erschlossene} Ergebnis nur so genau ist wie das Mittel aus 50 direkten Vergleichen. Nur wenn $q=\infty$, also $C-B$ absolut genau ist, wird das \textls{erschlossene} Ergebnis für $C-A$ ebenso genau wie das verwendete $B-A$, wie es auch aus der Natur der Sache folgt. \textsc{Littrows} falsche Formel würde für $q=\infty$ das absurde Ergebnis $4p$ ergeben. --- Ich vermute, dass \textsc{Littrow} dadurch irregeleitet wurde, dass der
